\section{Conclusion and Future Work}
\label{sec:conclusion}

This study evaluated the efficacy of HyDE against a standard RAG baseline within a highly specialized medical domain. Our findings demonstrate that while standard semantic search struggles with the vocabulary mismatch inherent in complex medical jargon and rare diseases, HyDE effectively bridges this gap. By generating an intermediate clinical document, HyDE significantly improved both the Hit Rate and Mean Reciprocal Rank (MRR) for these challenging query categories, equipping the ReAct agent with highly accurate context.

However, our results also highlight that HyDE is not universally necessary. For lexically dense queries---such as multi-sentence symptom descriptions---and common colloquialisms, standard RAG performed equally well as HyDE. This underscores the robust internal knowledge of the base Qwen2.5 7B model, which successfully mapped common terminology to clinical targets without the need for an intermediate generation step. Given the nearly 1.5x increase in generation time and higher token consumption introduced by HyDE, our evaluation suggests that an optimal production system should employ dynamic query routing. Such a system would utilize standard RAG for lexically rich or common queries, reserving the computationally expensive HyDE pipeline exclusively for sparse, jargon-heavy prompts.

While the absolute performance gains of HyDE may appear incremental in certain categories, in the high-stakes domain of healthcare, marginal improvements in retrieval accuracy are critically important. The computational trade-offs are heavily outweighed by the necessity of providing the reasoning agent with the most comprehensive and factually accurate context possible to prevent diagnostic hallucination. We do, however, acknowledge the limitations of evaluating this architecture purely through automated URL matching on a synthetic dataset. To truly gauge clinical viability and safety, the system must undergo rigorous qualitative stress-testing and human-in-the-loop evaluation by licensed medical professionals.

Looking forward, the structural richness of the Medline dataset offers a promising avenue for advancing this retrieval architecture. During our data ingestion phase, we systematically stripped relational XML tags, such as \texttt{<related\_pages>}, to optimize the raw text for dense vector embeddings. However, these discarded tags represent an explicit, human-curated web of medical relationships. Future work will focus on transitioning from standard vector search to a Graph RAG architecture. By embedding the medical documents as nodes and utilizing the \texttt{<related\_pages>} tags as relational edges, an autonomous agent could dynamically traverse this knowledge graph. This would enable multi-hop reasoning, allowing the agent to gather comprehensive, interconnected context for complex diagnostic queries that standard vector similarity currently fails to capture.